% =============================================================================
%  Lab 3 - Multiplexers
%  Seven sections. Delete the guidance comments as you fill each one in.
% =============================================================================

\chapter{Lab 3 --- Multiplexers}

\section{Objective}
% Two or three sentences, your own words.

\section{Specification}

\begin{table}[htbp]
\centering
\caption{Compressed table (Type C) for \texttt{Mux3}. Rows name a \emph{source},
not a value.}
\label{tab:lab03-truth}
\begin{tabular}{clc||ll}
\toprule
\# & \texttt{Sel1} & \texttt{Sel0} & $Y$ & note \\
\midrule
0 & 0 & 0 & \textbf{A} & \\
1 & 0 & 1 & \textbf{B} & \\
2 & 1 & 0 & \textbf{C} & \\
3 & 1 & 1 & \textbf{0} & defined default \\
\bottomrule
\end{tabular}
\end{table}

% ONE SENTENCE, required: why is this table Type C and not Type E? Give the
% number of rows an exhaustive table would need.

\subsection*{Boolean derivation}

\begin{align*}
\text{(1) canonical} \quad & Y_i = \overline{S_1}\,\overline{S_0}\,A_i + \overline{S_1}S_0\,B_i + S_1\overline{S_0}\,C_i + S_1 S_0 \cdot 0 \\
\text{(2) factored}  \quad & Y_i = \overline{S_1}\left( \overline{S_0}A_i + S_0 B_i \right) + S_1\overline{S_0}\,C_i
   \qquad \text{[\emph{name the law here}]} \\
\text{(3) VHDL} \quad & \text{see Listing~\ref{lst:code/lab03/Mux3.vhd}}
\end{align*}

% ONE SENTENCE, required: what familiar component is the bracketed
% sub-expression in line (2), and what does that tell you about the relationship
% between factoring an equation and decomposing a circuit?

\section{Design}
% Answer both:
%   (a) Which coding style did you use for Mux2 - when/else, with/select, or
%       case - and WHY? This is the graded question, not the syntax.
%   (b) Why does using 'case' with 'when others' make the Example 3 latch
%       impossible, while an if/elsif chain does not?

\section{Implementation}

\vhdlfile{code/lab03/Mux3.vhd}{\texttt{Mux3} --- generic-width 3-to-1 multiplexer}

\vhdlfile{code/lab03/CPUMuxes.vhd}{\texttt{CPUMuxes} --- the four multiplexers instantiated at CPU widths}

\section{Verification}
% Vectors, then the coverage argument. For a Type C table the argument has a
% fixed shape: exhaustive on the SELECTOR, distinct patterns on the DATA.

\begin{table}[htbp]
\centering
\caption{Test vectors. Note the data patterns are all distinct.}
\label{tab:lab03-vectors}
\begin{tabular}{cllll}
\toprule
\# & \texttt{Sel} & A / B / C & expected $Y$ & what wrong design this rules out \\
\midrule
1 & & & & \\
2 & & & & \\
3 & & & & \\
\bottomrule
\end{tabular}
\end{table}

% Answer both, one sentence each:
%   (a) Which vector catches a "multiplexer" that ignores its selector and
%       always returns A?
%   (b) Which vector catches a missing row 3, i.e. an inferred latch?

\section{Results}
% Evidence: E1, E2, E7.

\figplaceholder{the annotated CPU block diagram, with your four multiplexers
circled and labelled with what each one selects --- make this full page}
% \resultfig{figs/lab03/cpu_diagram_annotated.png}
%   {The four multiplexers built in this lab, shown in position on the datapath:
%    DselMUX selects ..., ALUinMUX selects ..., DinMUX selects ...}
%   {fig:lab03-diagram}

\figplaceholder{the RTL Viewer for \texttt{MUX3x32}}
% Required sentence: state that there are NO storage elements, and why that
% matters for a design that is supposed to be combinational.

\figplaceholder{the Questa waveform covering all four selector values}

% E1: state the error and warning counts. Say explicitly that there is no
%     inferred-latch warning.
% E4: state the register count for MUX3x32. It must be ZERO. Say why a non-zero
%     count would mean the design is wrong, not merely expensive.

\section{Discussion}
% REQUIRED. One bug, how you found it, what it cost.
% Good candidate: if you hit the inferred latch, describe how you noticed - the
% warning, the RTL Viewer, or the waveform holding a stale value - and which of
% those three told you fastest.
